\section*{Chapter \ref{ch:m2:money-markets} --- Money Markets}

\begin{solution}{exo:m2:money-markets:1}
$100\,(1 - 0.0375 \times 182/360) = 98.104167$ per 100.
\end{solution}

\begin{solution}{exo:m2:money-markets:2}
\omterm{def:m2:money-markets:yields}{Discount yield} $(100 - 99.04)/100 \times 360/91 = 3.7978\%$; \omterm{def:m2:money-markets:yields}{money-market yield}
$0.96/99.04 \times 360/91 = 3.8346\%$; \omterm{def:m2:money-markets:yields}{bond-equivalent yield} $0.96/99.04
\times 365/91 = 3.8879\%$.
\end{solution}

\begin{solution}{exo:m2:money-markets:3}
$10\,000\,000 \times (1 + 0.04 \times 90/360) = \text{USD}~10\,100\,000$. The CD's
rate is simple interest on the amount invested, actual/360: it is already a
\omterm{def:m2:money-markets:yields}{money-market yield}, directly comparable with the $m$ of a bill.
\end{solution}

\begin{solution}{exo:m2:money-markets:4}
By \cref{prop:m2:money-markets:order}, $i = \frac{y}{360}\,\frac{d}{1 -
dt/360}$, a product of two factors above one times $d$. The gap $m - d =
\frac{d^2 t/360}{1 - dt/360}$ grows with $t$: at the same discount rate the
26-week bill's yields exceed its discount rate by more than the 4-week's,
because the discount is taken on face value and the longer bill's price is
further below it.
\end{solution}

\begin{solution}{exo:m2:money-markets:5}
$P = 100\,(1 - 0.0765 \times 364/360) = 92.265000$. With $y = 365$: $a =
364/730 - 0.25 = 0.248630$, $b = 364/365 = 0.997260$, $c = (92.265 -
100)/92.265 = -0.083835$, and $i = (-b + \sqrt{b^2 - 4ac})/(2a) = 8.237\%$. A
coupon security of the same maturity would pay a coupon after six months,
which its holder could reinvest; the bill pays nothing until maturity. The
formula credits the bill with that reinvestment so that its yield compares
with a semiannual bond yield.
\end{solution}

\begin{solution}{exo:m2:money-markets:6}
Net redemptions are $1.1/20 = 5.5\%$ of net assets, above the 5\% threshold: an
institutional prime fund must charge the mandatory liquidity fee to the
redeeming investors. It pays the redemptions from its daily liquid assets, at
least 25\% of the fund. A government fund has no mandatory fee; it pays from
its liquid assets, which for a fund of bills and repo are most of the fund.
\end{solution}

\begin{solution}{exo:m2:money-markets:7}
4 weeks 3.8642\%, 6 weeks 3.8699\%, 8 weeks 3.8757\%, 13 weeks 3.8901\%, 17
weeks 3.9018\%, 26 weeks 3.9282\%, 52 weeks 3.9675\%. A flat discount curve
is an \emph{upward-sloping} yield curve: the conversion adds more at longer
terms (\cref{exo:m2:money-markets:4}), and beyond half a year the
semiannual formula takes over. Ten basis points of slope come from the
convention alone.
\end{solution}

\begin{solution}{exo:m2:money-markets:8}
In bond-equivalent terms the 13-week bill yields 3.8901\% and the 52-week
3.9567\%: the curve slopes \emph{up} by 6.7 basis points. Discount rates of
different maturities are not comparable yields. To read the market's
expectations one must also allow for term premia and for the supply of each
maturity, but first one must compare like with like.
\end{solution}

\begin{solution}{pb:m2:money-markets:1}
\textbf{1.} Seven days; the 31 December placement runs from Thursday to Monday,
four days.
\textbf{2.} $500 \times 10^6 \times 0.0395 \times 7/360 = \text{USD}~384\,028$.
\textbf{3.} $3.95\% \times 365/360 = 4.0049\%$.
\textbf{4.} $500 \times 10^6 \times [(1 + 0.0375/360)^3(1 + 0.0375 \times 4/360) -
1] = \text{USD}~364\,665$.
\textbf{5.} The deposit: the treasurer lends unsecured to the bank. The
facility: none, but the treasurer cannot use it; only eligible funds and
dealers can.
\textbf{6.} $500 \times 10^6 \times [(1 + 0.0388/360)^3 - 1] = \text{USD}~161\,684$.
\textbf{7.} $r_T = [(1 + 0.0395 \times 7/360)/(1 + 0.0388/360)^3 - 1] \times
360/4 = 4.0009\%$.
\textbf{8.} $4.0009 - 3.88 = 12.09$ basis points.
\textbf{9.} $500 \times 10^6 \times 0.001209 \times 4/360 = \text{USD}~6\,716$.
\textbf{10.} Its positions (securities it holds, loans it has made) must be
funded on that night whatever happens; the lenders who can still take the
risk and the balance sheet charge for them. What it cannot do cheaply on
one night is sell its assets.
\textbf{11.} On the reporting date dealers and banks took less of the funds'
cash; the funds placed it with the Fed, 245 billion more in one day, and took
it back on 3 January.
\textbf{12.} $500 \times 10^6 \times (0.040009 - 0.0375) \times 4/360 =
\text{USD}~13\,938$.
\textbf{13.} $[(1 + 0.0388/360)^{10}(1 + 0.040009 \times 4/360) - 1] \times
360/14 = 3.9171\%$.
\textbf{14.} The same four expensive days are averaged over fourteen instead of
seven: the turn adds 3.4 basis points to the two-week rate against 6.9 to the
one-week.
\textbf{15.} Not a coincidence in kind: eligible banks can borrow against
Treasuries from the Fed at 4.00\%, so a secured rate much above it would
be arbitraged by those with collateral and balance sheet to spare. An
unsecured deposit can exceed it, because the facility needs collateral and,
like any borrowing, enlarges the borrower's year-end balance sheet. A turn
far above would signal stress, or reluctance to be seen using the facility.
\textbf{16.} A one-month deposit from 1 December spans the year-end: its
rate contains the turn, diluted over 31 days, from the day it is quoted.
\textbf{17.} Measuring balance sheets as averages over the period rather than
on its last day; reserves so abundant that no bank needs to borrow over the
turn; central-bank facilities that take the cash without cost to anyone's
balance sheet (which the \omterm{def:m2:money-markets:rrp}{reverse repo facility} already does for funds).
\textbf{18.} Gains: no bank credit risk, a liquid asset. Risks and costs:
bills maturing just after the year-end are in demand and usually yield
\emph{less} than the deposit; she must buy and settle the bill, and she
loses the \omterm{def:m2:money-markets:turn}{turn premium} that the bank pays.
\textbf{19.} Named result: \emph{the implied year-end turn} is 4.0009\%, a
premium of 12.1 basis points over ordinary nights, for the four days from 31
December.
\textbf{20.} Because it is the night on which the size of balance sheets is
measured, and so the night on which balance sheet is scarce.
\end{solution}

\begin{solution}{iq:m2:money-markets:1}
The discount rate expresses the discount on the face value; the investor pays
less than face, so the return on what she actually pays is higher, and the
Treasury's \omterm{def:m2:money-markets:yields}{bond-equivalent yield} also uses a 365-day year instead of 360.
Compare with a two-year note on the \omterm{def:m2:money-markets:yields}{bond-equivalent yield} (for bills beyond
half a year, the semiannual formula), since note yields compound
semiannually on a 365-day basis.
\iqlookfor{face value against price paid, and 360 against 365.}
\end{solution}

\begin{solution}{iq:m2:money-markets:2}
The rate for a night on which balance sheets are measured (quarter-ends,
especially the year-end) jumps, because banks and dealers shrink repo and
deposits for that night; lenders who still offer balance sheet charge for it,
and cash that finds no borrower goes to central-bank facilities. The premium
appears in every term rate that spans the date, diluted by its length.
\iqlookfor{the balance-sheet measurement as cause, and the term-rate
dilution.}
\end{solution}

\begin{solution}{iq:m2:money-markets:3}
Prime funds promised a stable dollar share and held paper. When one fund's
shares fell below a dollar because of one issuer's default, investors in
every prime fund had a reason to redeem first, before losses were shared.
Funds met redemptions by not rolling paper and by selling it; issuers who
relied on rolling paper could not, and the market in new paper stopped. It
took a Treasury guarantee of fund shares and central-bank facilities to
restart it.
\iqlookfor{the first-mover incentive of a stable-value claim on risky assets,
and the rollover dependence of issuers.}
\end{solution}

\begin{solution}{iq:m2:money-markets:4}
The Fed is a riskless counterparty with unlimited capacity at a fixed rate
(up to the counterparty limit); dealers may not take all the cash, or only at
a lower rate on reporting dates; a fund may have counterparty limits with
each dealer, or want to keep the cash available late in the day when the
private market has closed. The 10 basis points are the price of those
differences.
\iqlookfor{capacity, counterparty limits and timing, not only the rate.}
\end{solution}

\begin{solution}{iq:m2:money-markets:5}
$r_T = [(1 + R N/B)/(1 + r_0/B)^{N_0} - 1] B/N_T$ with $N$ the term's days,
$N_0$ the ordinary days, $N_T$ the days the turn fixing covers (four if the
year-end falls on a Thursday before a holiday). Assumptions: ordinary nights
earn $r_0$ each, compounded daily; no other turn in the period; the term and
overnight quotes carry the same credit; the basis of 360 days.
\iqlookfor{the weights of the turn fixing, and naming the assumptions.}
\end{solution}

\begin{solution}{iq:m2:money-markets:6}
Start from the policy rate path: overnight-index swaps and overnight-rate
futures give compounded overnight rates to each date. Bills add the
government-supply and safety component; repo adds the collateral component;
build a curve per instrument class and model their spreads. Put turns in
explicitly: a jump in the forward overnight rate for each quarter-end and
year-end night, sized from quotes that straddle them, rather than letting
the interpolation smear them over neighbouring dates. Check by repricing
every input and by the forward rate's shape around each turn.
\iqlookfor{an explicit turn model, and separate curves for instruments with
different credit and collateral.}
\end{solution}
